% Extracción quirúrgica del propietario V2 05a: líneas 512--523.
% El resto permanece en este archivo como procedencia, pero no se publica.
\iffalse
\chapter{Macrostructure spectrale--polaire interne du continu}
\label{chap:continuo-macro-sp-interna}

\section{Portée strictement interne}

La construction de ce chapitre se termine avant toute reconnaissance
externe. Ses seules données sont le support APP à deux feuillets, l'orientation
TRIT, les extensions et la retenue TPK, le système projectif d'histoires, la
mesure cylindrique et le ledger signé. Le résultat est une macrostructure
graduée avec transport, expression de changement de feuillet, lecteur signé et
télescopage non autonome. Aucune fonction extérieure n'est utilisée pour définir
l'une quelconque de ses applications.

\section{Domaine survivant de la branche}

Pour \(x_k\in\mathcal A_k\), soit
\begin{equation}
\label{eq:continuo-sp-msect-n}
 \operatorname{Msect}^{\mathrm{sp},(N)}_k(x_k)
 =\left\{(x_j)_{j=k}^{N}:
 \tau_{j+1,j}x_{j+1}=x_j,
 \ (x_j,x_{j+9})\in\Gamma_{9,j},
 \ \text{et le faisceau est conservé }\{\Sigma,\Pi\}\right\}.
\end{equation}
On définit
\begin{equation}
\label{eq:continuo-sp-dominio-superviviente}
 \mathcal C_k^{\mathrm{sp}}
 =\left\{x_k:
 \operatorname{Msect}^{\mathrm{sp},(N)}_k(x_k)\neq\varnothing
 \text{ pour tout }N\geq k\right\},
 \qquad
 \mathcal C_{\rm cont}^{\mathrm{sp}}
 =\varprojlim_k\mathcal C_k^{\mathrm{sp}}.
\end{equation}
Le seuil TRIT n'est pas exclu de ce domaine. Le faisceau conserve les deux feuillets
arithmétiques même si une carte visible n'en publie aucun lors d'un
événement concret.

\section{Paquet de treize coordonnées}

Pour \(x_k\in\mathcal C_k^{\mathrm{sp}}\), le paquet dépendant est
\begin{equation}
\label{eq:continuo-sp-d13}
\begin{aligned}
 D_{\mathrm{sp},k}(x_k)=\bigl(&
 \mathcal F_{\mathrm{sp},k}(x_k),
 \operatorname{Ext}_{\Gamma,\mathrm{sp},k}(x_k),
 \operatorname{Rel}_{\mathrm{sp},k}(x_k),
 \mathcal N_{\mathrm{sp},k}(x_k),
 \operatorname{or}_{\mathrm{sp},k}(x_k),\\
 &\operatorname{Carr}_{\mathrm{sp},k}(x_k),
 \operatorname{Mem}_{\mathrm{sp},k}(x_k),
 \partial_{\mathrm{sp},k}(x_k),
 \gamma_{\mathrm{sp},k}(x_k),
 \operatorname{Msect}_{\mathrm{sp},k}(x_k),\\
 &\operatorname{Surv}_{\mathrm{sp},k}(x_k),
 \operatorname{Term}_{\mathrm{sp},k}(x_k),
 \mathcal L^{\rm sgn}_{\mathrm{sp},k}(x_k)\bigr).
\end{aligned}
\end{equation}
Ses composantes sont déterminées corrélativement comme suit.

La fibre \(\mathcal F_{\mathrm{sp},k}\) contient l'état TPK et le faisceau
ordonné \(\{\Sigma,\Pi\}\). L'extension est la restriction complète de
\(\operatorname{Ext}_\Gamma\) aux histoires qui satisfont
\eqref{eq:continuo-sp-dominio-superviviente}. Les relations contiennent les
lois d'identité et de composition, le relèvement résidu--quotient, le cocycle de
retenue, la compatibilité de frontière et la naturalité de la troncature. La
coordonnée numérique contient profondeur, phase, résidus, quotients et signature
interne. L'orientation stocke séparément feuillet et orientation TRIT. La
mémoire contient le quotient, le ledger et l'incrément de tour. La
frontière et le chemin restent non contractés. La multisection est le système
complet des prolongements de \eqref{eq:continuo-sp-msect-n}. La
survie est l'appartenance stable de
\eqref{eq:continuo-sp-dominio-superviviente}. Le terminal et le lecteur sont
construits dans les sections suivantes.

Les troncatures de toutes les coordonnées forment une application
\(u^{\mathrm{sp}}_{\ell,k}\) qui satisfait
\begin{equation}
\label{eq:continuo-sp-d-naturalidad}
 u^{\mathrm{sp}}_{\ell,k}D_{\mathrm{sp},\ell}
 =D_{\mathrm{sp},k}\tau_{\ell,k}.
\end{equation}

\section{Graphe de la restriction et section dépendante}

On définit
\begin{equation}
\label{eq:continuo-sp-g-graph}
 \mathcal G_{\mathrm{sp},k}
 =\operatorname{Graph}(D_{\mathrm{sp},k})
 =\{(x_k,D_{\mathrm{sp},k}x_k):
     x_k\in\mathcal C_k^{\mathrm{sp}}\}.
\end{equation}
La restriction totale est
\begin{equation}
\label{eq:continuo-sp-res-graph}
 \operatorname{Res}_{\mathrm{sp}}=s_{\mathrm{sp}}:
 \mathcal C_{\rm cont}^{\mathrm{sp}}
 \longrightarrow\mathcal G_{\mathrm{sp}},
 \qquad
 s_{\mathrm{sp}}(x)=(x,D_{\mathrm{sp}}x).
\end{equation}
La projection \(\pi_1(x,d)=x\) vérifie
\begin{equation}
\label{eq:continuo-sp-res-inversa}
 \pi_1s_{\mathrm{sp}}=I,
 \qquad
 s_{\mathrm{sp}}\pi_1=I_{\mathcal G_{\mathrm{sp}}}.
\end{equation}

Soit \(r_{\mathrm{sp},k}(x_k)\) la multisection des feuillets et des chemins qui appartient
déjà à \(D_{\mathrm{sp},k}(x_k)\). On définit
\begin{equation}
\label{eq:continuo-sp-x-dependiente}
 X_{\chi,\mathrm{sp},k}
 =\{(r_{\mathrm{sp},k}x_k,x_k,D_{\mathrm{sp},k}x_k):
     x_k\in\mathcal C_k^{\mathrm{sp}}\},
\end{equation}
\begin{equation}
\label{eq:continuo-sp-pr-x}
 \operatorname{pr}_{X,\mathrm{sp}}(x,D_{\mathrm{sp}}x)
 =(r_{\mathrm{sp}}x,x,D_{\mathrm{sp}}x).
\end{equation}
Oublier la première coordonnée de
\eqref{eq:continuo-sp-x-dependiente} est l'inverse de
\(\operatorname{pr}_{X,\mathrm{sp}}\).

\section{Faisceau de feuillets et mesure projective}

Soit
\begin{equation}
\label{eq:continuo-sp-haz-hojas}
 \widetilde{\mathcal C}_k^{\mathrm{sp}}
 =\{(x_k,h):x_k\in\mathcal C_k^{\mathrm{sp}},
              h\in\{\Sigma,\Pi\}\}.
\end{equation}
La mesure projective du continu induit des probabilités
\(\widetilde\mu_k\) sur ce faisceau. On se restreint au support positif ; pour
\(z\) dans ce support,
\begin{equation}
\label{eq:continuo-sp-consistencia-medida}
 \widetilde\mu_k(z)
 =\sum_{\widetilde\tau_{k+1,k}(z')=z}
  \widetilde\mu_{k+1}(z').
\end{equation}
La linéarisation de niveau est
\begin{equation}
\label{eq:continuo-sp-hilbert-nivel}
 \mathcal H_k^{\mathrm{sp}}
 =\ell^2(\operatorname{supp}\widetilde\mu_k),
\end{equation}
de base orthonormale \(e_z\).

\section{Transport de toutes les extensions et retenue}

Pour \(z'\) qui prolonge \(z\), nous définissons le poids conditionnel positif
\begin{equation}
\label{eq:continuo-sp-peso-condicional}
 w_k(z'\mid z)
 =\left(\frac{\widetilde\mu_{k+1}(z')}
               {\widetilde\mu_k(z)}\right)^{1/2}.
\end{equation}
Le transport est
\begin{equation}
\label{eq:continuo-sp-u-accion}
 U_ke_z
 =\sum_{\substack{z':\widetilde\tau_{k+1,k}(z')=z\\
            (z,z')\in\operatorname{Ext}_{\Gamma,\mathrm{sp},k}}}
 w_k(z'\mid z)e_{z'}.
\end{equation}
La somme contient toutes les extensions admissibles. Elle ne dépend pas d'un sélecteur.

Par \eqref{eq:continuo-sp-consistencia-medida},
\[
 \sum_{z':\widetilde\tau(z')=z}|w_k(z'\mid z)|^2=1.
\]
Les ensembles de descendants de deux parents distincts sont disjoints. Par
conséquent,
\begin{equation}
\label{eq:continuo-sp-u-isometria}
 U_k^*U_k=I_{\mathcal H_k^{\mathrm{sp}}}.
\end{equation}

Chaque base \(e_{z'}\) inclut la retenue de l'extension correspondante. Pour
des chemins composables, on conserve
\begin{align}
 \kappa(\gamma_2\gamma_1)
 &=\kappa(\gamma_2)+
   \mathsf C(\gamma_2)\kappa(\gamma_1),
 \label{eq:continuo-sp-carry-nativo}\\
 \operatorname{Carr}(\gamma_2\gamma_1)
 &=\operatorname{Carr}(\gamma_2)H(\gamma_1)
  +Y(\gamma_2)\operatorname{Carr}(\gamma_1).
 \label{eq:continuo-sp-carry-operador}
\end{align}
Les poids conditionnels se multiplient le long de l'unique chaîne de
troncatures d'un descendant. Les équations
\eqref{eq:continuo-sp-carry-nativo}--
\eqref{eq:continuo-sp-carry-operador} garantissent que les étiquettes de retenue
coïncident. Par conséquent,
\begin{equation}
\label{eq:continuo-sp-u-composicion}
 U_{\ell,k}=U_{\ell-1}\cdots U_k.
\end{equation}

\section{Graduation, projecteurs et décomposition du transport}

La graduation de feuillet est l'involution
\begin{equation}
\label{eq:continuo-sp-sigma}
 \sigma_ke_{(x,\Sigma)}=+e_{(x,\Sigma)},
 \qquad
 \sigma_ke_{(x,\Pi)}=-e_{(x,\Pi)}.
\end{equation}
Ainsi,
\[
 \sigma_k^*=\sigma_k,
 \qquad
 \sigma_k^2=I.
\]
Les projecteurs internes sont
\begin{equation}
\label{eq:continuo-sp-p-q}
 P_k=\frac{I+\sigma_k}{2},
 \qquad
 Q_k=\frac{I-\sigma_k}{2}.
\end{equation}
Elles satisfont
\begin{equation}
\label{eq:continuo-sp-p-q-relaciones}
 P_k^2=P_k,
 \quad Q_k^2=Q_k,
 \quad P_kQ_k=0,
 \quad P_k+Q_k=I.
\end{equation}

Nous séparons la partie du transport qui conserve le feuillet de celle qui le change :
\begin{align}
 A_k&=P_{k+1}U_kP_k+Q_{k+1}U_kQ_k,
 \label{eq:continuo-sp-a}\\
 \eta_k&=Q_{k+1}U_kP_k+P_{k+1}U_kQ_k.
 \label{eq:continuo-sp-eta}
\end{align}
Alors
\begin{equation}
\label{eq:continuo-sp-u-a-eta}
 U_k=A_k+\eta_k,
\end{equation}
et la graduation vérifie
\begin{equation}
\label{eq:continuo-sp-paridad-a-eta}
 \sigma_{k+1}A_k=A_k\sigma_k,
 \qquad
 \sigma_{k+1}\eta_k=-\eta_k\sigma_k.
\end{equation}

\begin{proposition}[Bilan local des feuillets]
\label{prop:continuo-sp-balance-local}
Pour chaque niveau,
\begin{equation}
\label{eq:continuo-sp-balance-local}
 A_k^*A_k+\eta_k^*\eta_k=I.
\end{equation}
\end{proposition}

\begin{proof}
L'orthogonalité de \(P_{k+1}\) et \(Q_{k+1}\) donne
\begin{align*}
 A_k^*A_k
 &=P_kU_k^*P_{k+1}U_kP_k
   +Q_kU_k^*Q_{k+1}U_kQ_k,\\
 \eta_k^*\eta_k
 &=P_kU_k^*Q_{k+1}U_kP_k
   +Q_kU_k^*P_{k+1}U_kQ_k.
\end{align*}
En sommant et en utilisant \(P_{k+1}+Q_{k+1}=I\),
\[
 A_k^*A_k+\eta_k^*\eta_k
 =P_kU_k^*U_kP_k+Q_kU_k^*U_kQ_k
 =P_k+Q_k=I.
\]
\end{proof}

L'opérateur \(\eta_k\) enregistre l'abandon du canal qui conserve le feuillet. Il
n'est pas la retenue entière : les coordonnées
\eqref{eq:continuo-sp-carry-nativo}--
\eqref{eq:continuo-sp-carry-operador} restent dans le paquet
\eqref{eq:continuo-sp-d13}.

\section{Lecteur signé avec ledger}

Pour \(v\in\mathcal H_k^{\mathrm{sp}}\), la polarisation signée est
\begin{equation}
\label{eq:continuo-sp-polarizacion-firmada}
 \operatorname{pol}_k(v)
 =\langle v,\sigma_kv\rangle
 =\|P_kv\|^2-\|Q_kv\|^2.
\end{equation}
Le lecteur spécialisé conserve en outre la source :
\begin{equation}
\label{eq:continuo-sp-lector-firmado}
 \mathcal L^{\rm sgn}_{\mathrm{sp},k}(x_k;v)
 =\bigl(
   \Lambda(\gamma_{x_k}),
   \Sigma_{\rm reg}(\Lambda(\gamma_{x_k})),
   \|P_kv\|^2,
   \|Q_kv\|^2,
   \operatorname{pol}_k(v)\bigr).
\end{equation}
L'orientation TRIT est conservée dans \(\Lambda(\gamma_{x_k})\) ; elle ne
s'identifie pas au signe de feuillet. À l'horizon \(N\), le registre garde la
famille des lecteurs pour tous les niveaux \(k\leq N\). Tronquer signifie
supprimer uniquement les entrées ultérieures.

\section{Produit non autonome et terminal}

On définit
\begin{equation}
\label{eq:continuo-sp-f-producto}
 F_{0,0}=I_{\mathcal H_0^{\mathrm{sp}}},
 \qquad
 F_{0,k}=A_{k-1}\cdots A_1A_0
 \quad(k\geq1).
\end{equation}
Par conséquent,
\begin{equation}
\label{eq:continuo-sp-f-recursion}
 F_{0,k+1}=A_kF_{0,k}.
\end{equation}
Le terme publié lors de l'abandon du canal à l'étape \(k\) et le terminal à
l'horizon \(N\) sont
\begin{equation}
\label{eq:continuo-sp-defecto-terminal}
 \mathcal D_k=F_{0,k}^*\eta_k^*\eta_kF_{0,k},
 \qquad
 \mathcal T_N=F_{0,N}^*F_{0,N}.
\end{equation}

\begin{theorem}[Télescopage spectral--polaire non autonome]
\label{thm:continuo-sp-telescopia-no-autonoma}
Pour tout horizon fini \(N\),
\begin{equation}
\label{eq:continuo-sp-telescopia-no-autonoma}
 \boxed{
 I=\sum_{k=0}^{N-1}
 F_{0,k}^*\eta_k^*\eta_kF_{0,k}
 +F_{0,N}^*F_{0,N}.}
\end{equation}
De plus,
\[
 0\preceq\mathcal T_{N+1}
 \preceq\mathcal T_N\preceq I,
\]
de sorte qu'il existe la limite forte positive
\[
 \mathcal T_\infty
 =\operatorname{s-lim}_{N\to\infty}\mathcal T_N
\]
et
\begin{equation}
\label{eq:continuo-sp-telescopia-infinita}
 I=\sum_{k=0}^{\infty}
 F_{0,k}^*\eta_k^*\eta_kF_{0,k}
 +\mathcal T_\infty
\end{equation}
en topologie forte. Le théorème n'affirme pas que
\(\mathcal T_\infty\) soit nul.
\end{theorem}

\begin{proof}
En multipliant \eqref{eq:continuo-sp-balance-local} par
\(F_{0,k}^*\) et \(F_{0,k}\), et en utilisant
\eqref{eq:continuo-sp-f-recursion}, on obtient
\[
 F_{0,k}^*F_{0,k}
 =F_{0,k+1}^*F_{0,k+1}
  +F_{0,k}^*\eta_k^*\eta_kF_{0,k}.
\]
La somme de \(k=0\) à \(N-1\) annule tous les termes
intermédiaires et laisse précisément
\eqref{eq:continuo-sp-telescopia-no-autonoma}. Comme
\(A_k^*A_k\preceq I\),
\[
 \mathcal T_{N+1}
 =F_{0,N}^*A_N^*A_NF_{0,N}
 \preceq F_{0,N}^*F_{0,N}=\mathcal T_N.
\]
La monotonie et la borne positive fournissent la limite forte. Le passage à la
limite dans l'identité finie donne
\eqref{eq:continuo-sp-telescopia-infinita}.
\end{proof}

\section{Assembleur comme graphe}

Pour
\[
 z_N=(r_{\mathrm{sp},N}x_N,x_N,D_{\mathrm{sp},N}x_N)
 \in X_{\chi,\mathrm{sp},N},
\]
on définit l'assembleur matériel
\begin{equation}
\label{eq:continuo-sp-a-accion}
\begin{aligned}
 a_{\mathrm{sp},N}(z_N)=\bigl(&
 (\mathcal H_k^{\mathrm{sp}},\sigma_k,P_k,Q_k)_{k\leq N},
 (U_k,A_k,\eta_k)_{k<N},\\
 &(F_{0,k},\mathcal D_k,\mathcal T_k,
   \mathcal L^{\rm sgn}_{\mathrm{sp},k})_{k\leq N}\bigr).
\end{aligned}
\end{equation}
La macrostructure est
\begin{equation}
\label{eq:continuo-sp-m-graph}
 \mathfrak M_{\mathrm{sp},N}
 =\operatorname{Graph}(a_{\mathrm{sp},N}),
 \qquad
 \mathcal A_{\mathrm{sp},N}(z_N)
 =(z_N,a_{\mathrm{sp},N}z_N).
\end{equation}
La projection sur \(z_N\) est l'inverse de
\(\mathcal A_{\mathrm{sp},N}\). La troncature supprime les entrées dont
l'indice dépasse le nouvel horizon et vérifie
\begin{equation}
\label{eq:continuo-sp-a-naturalidad}
 m^{\mathrm{sp}}_{\ell,k}a_{\mathrm{sp},\ell}
 =a_{\mathrm{sp},k}\xi^{\mathrm{sp}}_{\ell,k}.
\end{equation}

\section{Publicateur comme graphe}

Le publicateur ne contracte pas la macrostructure en un scalaire. Il renvoie le
certificat complet
\begin{equation}
\label{eq:continuo-sp-p-accion}
 p_{\mathrm{sp},N}(m_N)
 =\left(
 \mathcal L^{\rm sgn}_{\mathrm{sp},k},
 \mathcal D_k,
 \mathcal T_k,
 I=\sum_{j=0}^{k-1}\mathcal D_j+\mathcal T_k
 \right)_{0\leq k\leq N}.
\end{equation}
On définit
\begin{equation}
\label{eq:continuo-sp-c-graph}
 \mathcal C_{\mathrm{sp},N}
 =\operatorname{Graph}(p_{\mathrm{sp},N}),
 \qquad
 \mathcal P_{\mathrm{sp},N}(m_N)
 =(m_N,p_{\mathrm{sp},N}m_N).
\end{equation}
La projection sur \(m_N\) est son inverse. Tronquer le certificat supprime seulement
les identités ultérieures, de sorte que
\begin{equation}
\label{eq:continuo-sp-p-naturalidad}
 c^{\mathrm{sp}}_{\ell,k}p_{\mathrm{sp},\ell}
 =p_{\mathrm{sp},k}m^{\mathrm{sp}}_{\ell,k}.
\end{equation}

\section{Théorème de la chaîne interne complète}

\begin{theorem}[Macrostructure spectrale--polaire engendrée par le continu]
\label{thm:continuo-sp-cadena-interna}
Les actions
\eqref{eq:continuo-sp-res-graph},
\eqref{eq:continuo-sp-pr-x},
\eqref{eq:continuo-sp-m-graph} et
\eqref{eq:continuo-sp-c-graph} forment la chaîne naturelle
\begin{equation}
\label{eq:continuo-sp-cadena-interna}
 \mathcal C_{\rm cont}^{\rm disc}
 \supset\mathcal C_{\rm cont}^{\mathrm{sp}}
 \xrightarrow[\cong]{\operatorname{Res}_{\mathrm{sp}}}
 \mathcal G_{\mathrm{sp}}
 \xrightarrow[\cong]{\operatorname{pr}_{X,\mathrm{sp}}}
 X_{\chi,\mathrm{sp}}
 \xrightarrow[\cong]{\mathcal A_{\mathrm{sp}}}
 \mathfrak M_{\mathrm{sp}}
 \xrightarrow[\cong]{\mathcal P_{\mathrm{sp}}}
 \mathcal C_{\mathrm{sp}}.
\end{equation}
Son contenu publié est le faisceau à deux feuillets, la graduation
\(\sigma_k\), les projecteurs \(P_k,Q_k\), le transport total \(U_k\),
la séparation \(U_k=A_k+\eta_k\), la retenue et la mémoire, le lecteur signé
et le télescopage
\eqref{eq:continuo-sp-telescopia-no-autonoma} avec terminal explicite. Chaque
objet conserve une projection canonique vers tous ses antécédents.
\end{theorem}

\begin{proof}
Les identités
\eqref{eq:continuo-sp-res-inversa} prouvent la première équivalence. La
définition dépendante \eqref{eq:continuo-sp-x-dependiente} prouve la
deuxième. Les graphes
\eqref{eq:continuo-sp-m-graph} et
\eqref{eq:continuo-sp-c-graph} prouvent les deux restantes. La naturalité des
quatre étapes est
\eqref{eq:continuo-sp-d-naturalidad},
\eqref{eq:continuo-sp-a-naturalidad} et
\eqref{eq:continuo-sp-p-naturalidad}, avec la naturalité de
\(r_{\mathrm{sp},k}\). Le contenu opératoriel est prouvé par
\eqref{eq:continuo-sp-u-isometria}, le
\cref{prop:continuo-sp-balance-local} et le
\cref{thm:continuo-sp-telescopia-no-autonoma}. Toutes les extensions apparaissent
dans \eqref{eq:continuo-sp-u-accion} ; aucun choix de
chemin n'intervient donc.
\end{proof}

\fi
\section{Retour de phase et évolution non réinitialisée de l'état commun}

La branche conserve expressément
\begin{equation}
\label{eq:continuo-sp-no-reset}
 g(k+9)=g(k),
 \qquad
 x_{k+9}\neq x_k\quad\text{en général}.
\end{equation}
Le second membre enregistre l'accroissement de mémoire et les changements possibles
de feuillet, de retenue, de cylindre et de frontière. Ni \(\sigma_k\), ni les projecteurs, ni
le lecteur signé ne contractent cette différence.

\iffalse
\section{Provenance}

Les deux feuillets APP, l'état TRIT enrichi, les fibres et extensions TPK,
la retenue, le retour sans réinitialisation, la survie et le ledger sont une
\texttt{ARCHITECTURE\_AUTORALE\_PRÉEXISTANTE}. L'obligation de conserver le
terminal fini est un \texttt{RÉSULTAT\_RÉCUPÉRÉ}. Le faisceau explicite de feuillets,
l'involution \(\sigma_k\), les projecteurs, la linéarisation qui somme toutes
les extensions, la décomposition \(U_k=A_k+\eta_k\), le produit non
autonome \(F_{0,k}\) et les graphes dépendants sont une
\texttt{FORMALISATION\_NOUVELLE}. Le
\cref{thm:continuo-sp-telescopia-no-autonoma,thm:continuo-sp-cadena-interna}
constitue le \texttt{CERTIFICAT\_NOUVEAU} de cette formalisation.

\section{Falsadores internos}

\begin{proposition}[Falsificateurs de la macrostructure]
\label{prop:continuo-sp-falsadores}
La conclusion du
\cref{thm:continuo-sp-cadena-interna} n'est pas transférable si l'une quelconque
des substitutions suivantes se produit :
\begin{enumerate}
 \item choisir un prolongement dans
       \eqref{eq:continuo-sp-u-accion} et effacer les autres ;
 \item remplacer \(\mathcal G_{\mathrm{sp}}\) par l'image extérieure de
       \(D_{\mathrm{sp}}\) ;
 \item effacer \((x,D_{\mathrm{sp}}x)\) lors de la formation de
       \(X_{\chi,\mathrm{sp}}\) ;
 \item identifier le feuillet à l'orientation TRIT ou effacer le seuil ;
 \item rompre la cohérence
       \eqref{eq:continuo-sp-consistencia-medida} ;
 \item effacer la retenue ou enfreindre
       \eqref{eq:continuo-sp-carry-nativo}--
       \eqref{eq:continuo-sp-carry-operador} ;
 \item identifier \(\eta_k\) à toute la retenue ;
 \item remplacer la famille \((A_k)_k\) par un seul opérateur sans démontrer
       la stationnarité ;
 \item supprimer \(F_{0,N}^*F_{0,N}\) à un horizon fini ;
 \item affirmer \(\mathcal T_\infty=0\) sans preuve supplémentaire ;
 \item reconstruire une histoire à partir de la polarisation signée ;
 \item utiliser \(g(k+9)=g(k)\) pour affirmer \(x_{k+9}=x_k\) ;
 \item introduire un domaine, une fonction ou un critère extérieur pour
       définir l'une quelconque des flèches internes ;
 \item introduire une sortie ultérieure comme antécédent de la branche.
\end{enumerate}
\end{proposition}
\fi
